\documentclass[10pt]{article}
\usepackage[margin=1in]{geometry}
\usepackage{graphicx}
\usepackage{float}
\usepackage{xcolor}
\usepackage{hyperref}
\hypersetup{colorlinks=true, linkcolor=black, urlcolor=blue}

\graphicspath{{../proto-figures/}}

\title{\textbf{meridian-proto} — descriptors to UI}
\author{Layouts derived from gRPC service descriptors, rendered via meridian-mui-kit}
\date{}

\begin{document}
\maketitle

\noindent
Nobody authored these views. Each one was derived from \emph{protobuf descriptors
alone}: \texttt{meridian-proto} read a service's \texttt{FileDescriptorSet},
recovered a resource model from its \texttt{google.api.*} annotations —
\texttt{resource}, \texttt{field\_behavior}, \texttt{resource\_reference} — and
projected that into a neutral \texttt{meridian.ui.v1.ViewDescriptor}.
\texttt{@savvifi/meridian-mui-kit} then rendered it in hermetic Chrome for Testing.

\medskip
\noindent
Three things in these figures come from annotations rather than guesswork, and are
what a schema-derived UI cannot otherwise get right:
\texttt{field\_behavior} decides which fields appear in which view, so
\texttt{OUTPUT\_ONLY} values show on a detail card and never in a form;
column order is ranked with the identifier first and enums high, because a state is
the highest-signal cell in an operations table; and list pagination is read from the
AIP-158 request shape rather than assumed.

\medskip
\noindent
Cell values are synthesised per column \emph{format} — this is a structure catalog,
not live data — but every label, column, panel and RPC binding shown is the real
projected output.

\input{../proto-figures/figures.tex}

\end{document}
